\IfFileExists{stacks-project.cls}{%
\documentclass{stacks-project}
}{%
\documentclass{amsart}
}

% For dealing with references we use the comment environment
\usepackage{verbatim}
\newenvironment{reference}{\comment}{\endcomment}
%\newenvironment{reference}{}{}
\newenvironment{slogan}{\comment}{\endcomment}
\newenvironment{history}{\comment}{\endcomment}

% For commutative diagrams we use Xy-pic
\usepackage[all]{xy}

% We use 2cell for 2-commutative diagrams.
\xyoption{2cell}
\UseAllTwocells

% We use multicol for the list of chapters between chapters
\usepackage{multicol}

% This is generally recommended for better output
\usepackage{lmodern}
\usepackage[T1]{fontenc}

% For cross-file-references
\usepackage{xr-hyper}

% Package for hypertext links:
\usepackage{hyperref}

% For any local file, say "hello.tex" you want to link to please
% use \externaldocument[hello-]{hello}
\externaldocument[introduction-]{introduction}
\externaldocument[conventions-]{conventions}
\externaldocument[sets-]{sets}
\externaldocument[categories-]{categories}
\externaldocument[topology-]{topology}
\externaldocument[sheaves-]{sheaves}
\externaldocument[sites-]{sites}
\externaldocument[stacks-]{stacks}
\externaldocument[fields-]{fields}
\externaldocument[algebra-]{algebra}
\externaldocument[brauer-]{brauer}
\externaldocument[homology-]{homology}
\externaldocument[derived-]{derived}
\externaldocument[simplicial-]{simplicial}
\externaldocument[more-algebra-]{more-algebra}
\externaldocument[smoothing-]{smoothing}
\externaldocument[modules-]{modules}
\externaldocument[sites-modules-]{sites-modules}
\externaldocument[injectives-]{injectives}
\externaldocument[cohomology-]{cohomology}
\externaldocument[sites-cohomology-]{sites-cohomology}
\externaldocument[dga-]{dga}
\externaldocument[dpa-]{dpa}
\externaldocument[sdga-]{sdga}
\externaldocument[hypercovering-]{hypercovering}
\externaldocument[schemes-]{schemes}
\externaldocument[constructions-]{constructions}
\externaldocument[properties-]{properties}
\externaldocument[morphisms-]{morphisms}
\externaldocument[coherent-]{coherent}
\externaldocument[divisors-]{divisors}
\externaldocument[limits-]{limits}
\externaldocument[varieties-]{varieties}
\externaldocument[topologies-]{topologies}
\externaldocument[descent-]{descent}
\externaldocument[perfect-]{perfect}
\externaldocument[more-morphisms-]{more-morphisms}
\externaldocument[flat-]{flat}
\externaldocument[groupoids-]{groupoids}
\externaldocument[more-groupoids-]{more-groupoids}
\externaldocument[etale-]{etale}
\externaldocument[chow-]{chow}
\externaldocument[intersection-]{intersection}
\externaldocument[pic-]{pic}
\externaldocument[weil-]{weil}
\externaldocument[adequate-]{adequate}
\externaldocument[dualizing-]{dualizing}
\externaldocument[duality-]{duality}
\externaldocument[discriminant-]{discriminant}
\externaldocument[derham-]{derham}
\externaldocument[local-cohomology-]{local-cohomology}
\externaldocument[algebraization-]{algebraization}
\externaldocument[curves-]{curves}
\externaldocument[resolve-]{resolve}
\externaldocument[models-]{models}
\externaldocument[functors-]{functors}
\externaldocument[equiv-]{equiv}
\externaldocument[pione-]{pione}
\externaldocument[etale-cohomology-]{etale-cohomology}
\externaldocument[proetale-]{proetale}
\externaldocument[relative-cycles-]{relative-cycles}
\externaldocument[more-etale-]{more-etale}
\externaldocument[trace-]{trace}
\externaldocument[crystalline-]{crystalline}
\externaldocument[spaces-]{spaces}
\externaldocument[spaces-properties-]{spaces-properties}
\externaldocument[spaces-morphisms-]{spaces-morphisms}
\externaldocument[decent-spaces-]{decent-spaces}
\externaldocument[spaces-cohomology-]{spaces-cohomology}
\externaldocument[spaces-limits-]{spaces-limits}
\externaldocument[spaces-divisors-]{spaces-divisors}
\externaldocument[spaces-over-fields-]{spaces-over-fields}
\externaldocument[spaces-topologies-]{spaces-topologies}
\externaldocument[spaces-descent-]{spaces-descent}
\externaldocument[spaces-perfect-]{spaces-perfect}
\externaldocument[spaces-more-morphisms-]{spaces-more-morphisms}
\externaldocument[spaces-flat-]{spaces-flat}
\externaldocument[spaces-groupoids-]{spaces-groupoids}
\externaldocument[spaces-more-groupoids-]{spaces-more-groupoids}
\externaldocument[bootstrap-]{bootstrap}
\externaldocument[spaces-pushouts-]{spaces-pushouts}
\externaldocument[spaces-chow-]{spaces-chow}
\externaldocument[groupoids-quotients-]{groupoids-quotients}
\externaldocument[spaces-more-cohomology-]{spaces-more-cohomology}
\externaldocument[spaces-simplicial-]{spaces-simplicial}
\externaldocument[spaces-duality-]{spaces-duality}
\externaldocument[formal-spaces-]{formal-spaces}
\externaldocument[restricted-]{restricted}
\externaldocument[spaces-resolve-]{spaces-resolve}
\externaldocument[formal-defos-]{formal-defos}
\externaldocument[defos-]{defos}
\externaldocument[cotangent-]{cotangent}
\externaldocument[examples-defos-]{examples-defos}
\externaldocument[algebraic-]{algebraic}
\externaldocument[examples-stacks-]{examples-stacks}
\externaldocument[stacks-sheaves-]{stacks-sheaves}
\externaldocument[criteria-]{criteria}
\externaldocument[artin-]{artin}
\externaldocument[quot-]{quot}
\externaldocument[stacks-properties-]{stacks-properties}
\externaldocument[stacks-morphisms-]{stacks-morphisms}
\externaldocument[stacks-limits-]{stacks-limits}
\externaldocument[stacks-cohomology-]{stacks-cohomology}
\externaldocument[stacks-perfect-]{stacks-perfect}
\externaldocument[stacks-introduction-]{stacks-introduction}
\externaldocument[stacks-more-morphisms-]{stacks-more-morphisms}
\externaldocument[stacks-geometry-]{stacks-geometry}
\externaldocument[moduli-]{moduli}
\externaldocument[moduli-curves-]{moduli-curves}
\externaldocument[examples-]{examples}
\externaldocument[exercises-]{exercises}
\externaldocument[guide-]{guide}
\externaldocument[desirables-]{desirables}
\externaldocument[coding-]{coding}
\externaldocument[obsolete-]{obsolete}
\externaldocument[fdl-]{fdl}
\externaldocument[index-]{index}

% Theorem environments.
%
\theoremstyle{plain}
\newtheorem{theorem}[subsection]{Theorem}
\newtheorem{proposition}[subsection]{Proposition}
\newtheorem{lemma}[subsection]{Lemma}

\theoremstyle{definition}
\newtheorem{definition}[subsection]{Definition}
\newtheorem{example}[subsection]{Example}
\newtheorem{exercise}[subsection]{Exercise}
\newtheorem{situation}[subsection]{Situation}

\theoremstyle{remark}
\newtheorem{remark}[subsection]{Remark}
\newtheorem{remarks}[subsection]{Remarks}

\numberwithin{equation}{subsection}

% Macros
%
\def\lim{\mathop{\mathrm{lim}}\nolimits}
\def\colim{\mathop{\mathrm{colim}}\nolimits}
\def\Spec{\mathop{\mathrm{Spec}}}
\def\Hom{\mathop{\mathrm{Hom}}\nolimits}
\def\Ext{\mathop{\mathrm{Ext}}\nolimits}
\def\SheafHom{\mathop{\mathcal{H}\!\mathit{om}}\nolimits}
\def\SheafExt{\mathop{\mathcal{E}\!\mathit{xt}}\nolimits}
\def\Sch{\mathit{Sch}}
\def\Mor{\mathop{\mathrm{Mor}}\nolimits}
\def\Ob{\mathop{\mathrm{Ob}}\nolimits}
\def\Sh{\mathop{\mathit{Sh}}\nolimits}
\def\NL{\mathop{N\!L}\nolimits}
\def\CH{\mathop{\mathrm{CH}}\nolimits}
\def\proetale{{pro\text{-}\acute{e}tale}}
\def\etale{{\acute{e}tale}}
\def\QCoh{\mathit{QCoh}}
\def\Ker{\mathop{\mathrm{Ker}}}
\def\Im{\mathop{\mathrm{Im}}}
\def\Coker{\mathop{\mathrm{Coker}}}
\def\Coim{\mathop{\mathrm{Coim}}}

% Boxtimes
%
\DeclareMathSymbol{\boxtimes}{\mathbin}{AMSa}{"02}

%
% Macros for moduli stacks/spaces
%
\def\QCohstack{\mathcal{QC}\!\mathit{oh}}
\def\Cohstack{\mathcal{C}\!\mathit{oh}}
\def\Spacesstack{\mathcal{S}\!\mathit{paces}}
\def\Quotfunctor{\mathrm{Quot}}
\def\Hilbfunctor{\mathrm{Hilb}}
\def\Curvesstack{\mathcal{C}\!\mathit{urves}}
\def\Polarizedstack{\mathcal{P}\!\mathit{olarized}}
\def\Complexesstack{\mathcal{C}\!\mathit{omplexes}}
% \Pic is the operator that assigns to X its picard group, usage \Pic(X)
% \Picardstack_{X/B} denotes the Picard stack of X over B
% \Picardfunctor_{X/B} denotes the Picard functor of X over B
\def\Pic{\mathop{\mathrm{Pic}}\nolimits}
\def\Picardstack{\mathcal{P}\!\mathit{ic}}
\def\Picardfunctor{\mathrm{Pic}}
\def\Deformationcategory{\mathcal{D}\!\mathit{ef}}
\title{The representing object of a truncated simplicial object is the limit over the truncated simplex category}
\author{The Stacks Project Authors}
\date{stacks-project:simplicial.tex:local-snapshot-20261012; GNU FDL 1.2 or later}
\begin{document}
\maketitle
\noindent\textit{Editorial note.} The counted statement is the author's own lemma environment labelled lemma-formula-limit (source lines 2553-2570); the proof environment begins two lines below it and is its own proof, closing at source line 2642.\par\medskip
\noindent\textit{Editorial note.} No \textbackslash{}ref or \textbackslash{}cite occurs anywhere inside the retained span, so the package declares neither a source pointer nor a bibliography entry; the truncated simplicial objects and the limit involved are developed earlier in the same chapter.\par\medskip
\noindent\textit{Editorial note.} The proof closes with the author's own sentence leaving the verification that the two constructions are mutually inverse bijections and functorial in T to the reader; the retained package reproduces that sentence and does not supply the verification.\par\medskip
\noindent\textit{Editorial note.} The retained preamble is the chapter's own minimal preamble (sources/stacks-simplicial/preamble.tex, taken with its real line count), a verbatim excerpt of the native Xy-pic setup, the native theorem declarations with the mathematical macros, and the project's own no-op reference environments; no macro was invented or redefined.\par\medskip
\noindent\textit{Editorial scope.} One original result of the Stacks Project's Simplicial Objects chapter, counted once: a representing object of the truncated tuple functor is the limit over the truncated simplex category.\par\medskip
\noindent\textit{Allowed standard prerequisites (editorial).}\begin{itemize}
\item n-truncated simplicial object. A sequence of objects U\_0, ..., U\_n with face and degeneracy maps satisfying the simplicial identities, regarded as a functor on the category of finite non-empty ordinals of size at most n+1.. The face maps and the identities are what the tuple condition in the statement and the verification in the proof use (source lines 2556-2563, 2620-2638).
\item Limit over the truncated simplex category. The limit of a truncated simplicial object over the opposite of the category of maps [k] -> [n+1] with k <= n, presented as the set of compatible families.. It is the right-hand side of the statement and the object compared with the representing tuple (source lines 2568, 2573-2580).
\item Representable functor. A contravariant functor to sets is representable by an object when there is a natural bijection between its values and the morphisms into that object.. The hypothesis of representability is what lets the proof compare the two functors pointwise on an arbitrary object T (source lines 2565-2566, 2581-2600).
\end{itemize}\medskip
\begin{lemma}
\label{lemma-formula-limit}
Let $n$ be an integer $\geq 1$.
Let $U$ be a $n$-truncated simplicial object of $\mathcal{C}$.
Consider the contravariant functor from $\mathcal{C}$ to
$\textit{Sets}$ which associates to an object $T$ the set
$$
\{ (f_0, \ldots, f_{n + 1}) \in \Mor_\mathcal{C}(T, U_n)
\mid
d^n_{j - 1} \circ f_i = d^n_i \circ f_j
\ \forall\ 0\leq i < j\leq n + 1\}
$$
If this functor is representable by some object $U_{n + 1}$
of $\mathcal{C}$, then
$$
U_{n + 1} = \lim_{(\Delta/[n + 1])_{\leq n}^{opp}} U(n)
$$
\end{lemma}

\begin{proof}
The limit, if it exists, represents the functor
that associates to an object $T$ the set
$$
\{
(f_\alpha)_{\alpha : [k] \to [n + 1], k \leq n}
\mid
f_{\alpha \circ \psi} = U(\psi) \circ f_\alpha\ \forall
\ \psi : [k'] \to [k], \alpha : [k] \to [n + 1]
\}.
$$
In fact we will show this functor is isomorphic to the
one displayed in the lemma. The map in one direction
is given by the rule
$$
(f_\alpha)_{\alpha}
\longmapsto
(f_{\delta^{n + 1}_0}, \ldots, f_{\delta^{n + 1}_{n + 1}}).
$$
This satisfies the conditions of the lemma because
$$
d^n_{j - 1} \circ f_{\delta^{n + 1}_i} =
f_{\delta^{n + 1}_i \circ \delta^n_{j - 1}} =
f_{\delta^{n + 1}_j \circ \delta^n_i} =
d^n_i \circ f_{\delta^{n + 1}_j}
$$
by the relations we recalled above the lemma. To construct a map
in the other direction we have to associate to a system
$(f_0, \ldots, f_{n + 1})$ as in the displayed formula
of the lemma a system of maps $f_\alpha$. Let $\alpha : [k] \to [n + 1]$
be given. Since $k \leq n$ the map $\alpha$ is not surjective.
Hence we can write $\alpha = \delta^{n + 1}_i \circ \psi$
for some $0 \leq i \leq n + 1$ and some
$\psi : [k] \to [n]$. We have no choice but to define
$$
f_\alpha = U(\psi) \circ f_i.
$$
Of course we have to check that this is independent of the
choice of the pair $(i, \psi)$. First, observe that given $i$
there is a unique $\psi$ which works. Second, suppose that $(j, \phi)$ is
another pair. Then $i \not = j$ and we may assume $i < j$. Since
both $i, j$ are not in the image of $\alpha$ we may actually
write $\alpha = \delta^{n + 1}_{i, j} \circ \xi$ and then
we see that $\psi = \delta^n_{j - 1} \circ \xi$ and
$\phi = \delta^n_i \circ \xi$. Thus
\begin{eqnarray*}
U(\psi) \circ f_i & = & U(\delta^n_{j - 1} \circ \xi) \circ f_i \\
& = & U(\xi) \circ d^n_{j - 1} \circ f_i \\
& = & U(\xi) \circ d^n_i \circ f_j \\
& = & U(\delta^n_i \circ \xi) \circ f_j \\
& = & U(\phi) \circ f_j
\end{eqnarray*}
as desired. We still have to verify that the maps
$f_\alpha$ so defined satisfy the rules of a system
of maps $(f_\alpha)_\alpha$. To see this suppose that
$\psi : [k'] \to [k]$, $\alpha : [k] \to [n + 1]$ with
$k, k' \leq n$. Set $\alpha' = \alpha \circ \psi$.
Choose $i$ not in the image of $\alpha$. Then clearly
$i$ is not in the image of $\alpha'$ also. Write
$\alpha = \delta^{n + 1}_i \circ \phi$ (we cannot use the letter $\psi$ here
because we've already used it). Then obviously
$\alpha' = \delta^{n + 1}_i \circ \phi \circ \psi$. By construction above
we then have
$$
U(\psi) \circ f_\alpha = U(\psi) \circ U(\phi) \circ f_i
= U(\phi \circ \psi) \circ f_i = f_{\alpha \circ \psi} = f_{\alpha'}
$$
as desired. We leave to the reader the pleasant task of verifying
that our constructions are mutually inverse bijections, and are
functorial in $T$.
\end{proof}

\end{document}
